\section{An equivalent conjectural package}\label{sec:program}

\subsection{The three central assertions}\label{sec:three-assertions}

Fix a product \(X=\prod_{i=1}^nC_i\), its points \(o_i\), and
\(\ell=\sum_i a_iH_i\), with \(a_i>0\). Pass to one field
\(F_0\) defining the maximal curve-visible projectors for all
odd subproducts. All height assertions use
\Cref{def:heights}.

\begin{conjecture}[Off-central vanishing]
\label[conjecture]{conj:offcentral}
For every subproduct \(Y_S\), every finite \(E/F_0\), and
every \(c\) with \(|S|\ne2c-1\),
\[
 U_S^c(E)=0.
\]
\end{conjecture}

\begin{conjecture}[Central algebraic Abel--Jacobi injectivity]
\label[conjecture]{conj:AJ}
For every odd subproduct \(|S|=2c-1\) and every finite \(E/F_0\),
the rationalized Abel--Jacobi map is injective on \(A_{S,E}\).
\end{conjecture}

\begin{conjecture}[Residual integral Northcott]
\label[conjecture]{conj:residual-northcott}
For every odd subproduct \(|S|=2c-1\), the quadratic form
\(q_{G,S}=(-1)^ch_{G,S}\) has bounded-degree Northcott on
the projected integral image \(\Lambda_{G,S}\).
\end{conjecture}

The real-positivity assertion in the central program is a
consequence of \Cref{conj:residual-northcott}, by
\Cref{lem:northcott-real}. It can still be studied first as the
weaker conjecture that \(q_{G,S}\) is real positive definite.
\Cref{ex:denominators} shows the additional integral issue in the
converse direction.

\begin{theorem}[Equivalence of the product and residual programs]
\label[theorem]{thm:program-equivalence}
For the fixed pointed product and the displayed polarization, the
following packages are equivalent.
\begin{enumerate}[label=(\roman*)]
\item Arithmetic hard Lefschetz holds in every degree and over
every finite extension of \(F_0\), and
\(q_{\ell,r}\) has bounded-degree Northcott on the primitive
integral image in every degree
\(1\le r\le\lfloor(n+1)/2\rfloor\).
\item \Cref{conj:offcentral,conj:AJ,conj:residual-northcott} hold
for all the subproducts of \(X\).
\end{enumerate}
Either package implies real primitive positivity and
nondegeneracy of the complementary height pairings over every
fixed number field.
\end{theorem}

\begin{proof}
From (ii), \Cref{thm:main-lefschetz} gives arithmetic hard
Lefschetz. \Cref{thm:central-northcott} gives central integral
Northcott, and \Cref{lem:northcott-real} gives real central
positivity. \Cref{cor:positivity-propagation} gives all the
primitive Northcott and positivity assertions, as well as
nondegeneracy.

Conversely, hard Lefschetz in (i) gives off-central vanishing by
\Cref{thm:main-lefschetz}. For an odd subproduct \(|S|=2c-1\),
inflate \(u\in U_S^c(E)\) by fundamental classes on all omitted
factors. In the notation of \MathSecref{sec:chow-blocks}, this is
the integral injection \(u\mapsto j_\varnothing u\) in degree
\(r=c\). With \(m=n-|S|\), it is primitive and satisfies
\[
 q_{\ell,c}(j_\varnothing u)
       =m!\prod_{i\notin S}a_i\,q_{S,c}(u).
\]
Thus primitive Northcott on \(X\) gives central integral
Northcott. \Cref{thm:central-northcott} supplies
\Cref{conj:AJ,conj:residual-northcott}. The final assertion now
follows from the first direction of the proof.
\end{proof}

This is a product-of-curves formulation of the arithmetic
Lefschetz and Hodge--Riemann program, with an integral finite-height
condition added. The underlying arithmetic standard conjectures
are those in \cite[Conjecture~2.3.1]{ZhangStandard}.

\subsection{Arithmetic detection and saturation}
\label{sec:detection-conjectures}

Fix one central residual group \(G_E\).

\begin{conjecture}[Residual rational detection]
\label[conjecture]{conj:detection}
For every finite \(E/F_0\) and every prime \(p\), the map
\[
 G_E\otimes_\Q\Q_p\longrightarrow H^1(E,V_p)
\]
of \Cref{prop:residual-regulator} is injective.
\end{conjecture}

This is a residual form of rational arithmetic cycle-class
injectivity, in the framework of \cite{Raskind,Kahn}. Together
with \Cref{thm:lattice-equivalence}, it reduces fixed-field
Northcott under real positivity to the following integral question.

\begin{conjecture}[Residual integral discreteness]
\label[conjecture]{conj:lattice}
For every finite \(E/F_0\), \(G_E\) is finite-dimensional and
there are a rational basis \(v_1,\ldots,v_s\) and an integer
\(N_E\ge1\) with
\[
 \Lambda_G(E)\subseteq N_E^{-1}\sum_i\Z v_i.
\]
\end{conjecture}

One route is to prove the common-lattice, almost-everywhere
saturation hypothesis of \Cref{thm:all-prime} for the arithmetic
regulators. A second route is a genuine positive height gap,
together with the finite Kummer quotients of
\Cref{thm:weak-mw-gap}. The local integral images of
\Cref{prop:local-lattice} are already controlled; their comparison
as the coefficient prime varies is the remaining step.

Integral detection must refer to the stated image in a rational
quotient. The unrestricted integral continuous cycle class map
has torsion counterexamples over number fields in codimension
three \cite[Theorems~1.3--1.4]{ScaviaSuzuki}. Those groups and
their torsion kernels are different from \(\Lambda_G(E)\).

\subsection{Varying fields and local metrics}\label{sec:next-targets}

\begin{conjecture}[Residual class-field ramification]
\label[conjecture]{conj:ramification}
For each \(B\ge0,D\ge1\), the minimal class fields
\(F_0(\xi)\) of \(\xi\in\Lambda_G(E)\), with
\([E:F_0]\le D\) and \(q_G(\xi)\le B\), are unramified
outside one finite set of places depending on \(Y,\Pi,B,D\).
\end{conjecture}

With fixed-field Northcott, this is exactly the remaining
bounded-degree condition by \Cref{thm:ramification}. A quantitative
version asks whether constants \(A_D,B_D\) in
\MathEqref{eq:discriminant-target} exist. The quantitative
inequality is recorded as a research question, rather than being
identified with the qualitative conjecture.

For the metric problem, the next comparison is local and precise.
Starting with a fixed semistable source curve, construct the
admissible correction to the model determinant and compare its
rigidified metric with the classical Poincaré metric over the
original valued field. \Cref{prop:good-reduction} is the smooth
model case; the complex-place comparison is
\Cref{thm:quillen-isometry,cor:cycle-determinant}. Applying these
comparisons to an actual-form non-Archimedean Green construction
requires its trace, covariance and normalization statements on
their full test spaces.

Boundary estimates in a varying family also require their own
geometric input. The smooth relative biextension comparison in
\cite[Theorem~10.9]{Nakayama} is unconditional. In that work,
the singular-fibre extension in Theorem~11.2 uses a modified cycle
satisfying a specified cohomological condition, and Lemma~11.1
obtains that modification from the Hodge conjecture. These are
the hypotheses to verify when importing boundary asymptotics
into the family-height program.
